\documentclass[11pt,a4paper]{article}
\usepackage{DDTA_DERMATRIAGE_GUIDED_REVIEW_STYLE_R1}

\hypersetup{
  pdftitle={DDTA R25 DermaTriage - Decision Review Reading Note 01},
  pdfauthor={Documentation-Driven Threat Analysis research project}
}

\pagestyle{fancy}
\fancyhf{}
\fancyhead[L]{\sffamily\small\color{DDTAGray}DDTA DermaTriage Decision Review}
\fancyhead[R]{\sffamily\small\color{DDTAGray}READING NOTE 01 - NOT REPOSITORY STATE}
\fancyfoot[L]{\sffamily\scriptsize\color{DDTAGray}Temporary reading artifact - Decision phase}
\fancyfoot[R]{\sffamily\small\color{DDTAGray}\thepage}
\renewcommand{\headrulewidth}{0.4pt}
\renewcommand{\footrulewidth}{0pt}

\begin{document}
\selectlanguage{italian}

\begin{titlepage}
\thispagestyle{empty}
\vspace*{9mm}
{\sffamily\bfseries\color{DDTANavy}\fontsize{15}{18}\selectfont Documentation-Driven Threat Analysis}\par
\vspace{5mm}
{\sffamily\bfseries\color{DDTANavy}\fontsize{28}{32}\selectfont DermaTriage Decision Review}\par
\vspace{2mm}
{\sffamily\color{DDTABlue}\fontsize{18}{22}\selectfont Reading Note 01 - applicazione esplicita dei gate a DEC-01}\par

\vspace{10mm}
\begin{tcolorbox}[enhanced,arc=3pt,boxrule=0pt,colback=DDTANavy,left=12pt,right=12pt,top=11pt,bottom=11pt]
{\color{white}\sffamily\bfseries\large Artifact temporaneo di lettura}\par
\vspace{2mm}
{\color{white!88}\small Questo documento non modifica R1, non costituisce R2 e non deve essere salvato nel repository. Serve a rendere visibili le domande dei gate prima di proseguire con le Decision successive.}
\end{tcolorbox}

\vspace{8mm}
\begin{tabularx}{\textwidth}{L{45mm}Y}
\toprule
\textbf{Repository baseline} & \texttt{bb9eb66aeb611c3f1630e39f5511690b7eb6287f} \\
\textbf{Working record} & R1 congelata e immutabile \\
\textbf{Metodo} & \texttt{DDTA\_DOCUMENTATION\_BA\_AUTHORING\_GUIDE\_R4} \\
\textbf{Parent element} & \texttt{MR-01 - Valutazione di triage del caso dermatologico} \\
\textbf{Elemento riesaminato} & \texttt{DEC-01 - Continuita' del triage in assenza di immagine} \\
\textbf{Base Analysis} & BLOCCATA fino al documentation gate \\
\textbf{Data} & 3 settembre 2026 \\
\bottomrule
\end{tabularx}

\vfill
\begin{ddtaprinciple}[title={Regola per questa fase}]
Per ogni Decision il working review record deve mostrare non solo l'esito del gate, ma anche la \textbf{domanda che il reviewer pone}, la risposta supportata dalla source evidence, l'eventuale gap, la disposition e la STOP decision.
\end{ddtaprinciple}
\end{titlepage}

\section{Perche' questa nota e' necessaria}

La R1 aveva gia' autorizzato una procedura esplicita per la prima Decision sotto MR-01:
\begin{enumerate}
  \item source evidence;
  \item candidate commitment;
  \item Decision gate;
  \item responsibility-boundary gate;
  \item Decision-vs-current-realization test;
  \item SELECTABLE / NECESSITY-CONSTRAINED / DEFAULT review;
  \item semantic-sufficiency check;
  \item explicit disposition e STOP decision.
\end{enumerate}

Nella prima esposizione di DEC-01 i criteri sono stati applicati, ma sono stati compressi in una tabella di esiti. Questo rendeva meno auditabile il percorso \emph{domanda -> evidence -> risposta -> disposition}.

\begin{ddtacore}[title={Classificazione della discrepanza}]
\textbf{Nessuna modifica metodologica proposta in questo momento.}
La regola era gia' presente nella guida e nella R1. La discrepanza viene classificata come \textbf{execution/presentation issue del working review}, non come method defect o metamodel pressure.
\end{ddtacore}

Resta aperta una osservazione metodologica: se il problema dovesse ripetersi, potra' essere utile rendere normativa nella guida la visibilita' delle \emph{gate questions}, non soltanto la registrazione dei gate applicati.

\section{Source evidence minima per DEC-01}

\subsection{Evidence E1 - percorso principale con immagine}
\textbf{Source:} \texttt{OR2\_Architecture\_Document.pdf}, sezione 2, pagina 1.

La fonte descrive un percorso principale nel quale il paziente fornisce un'immagine della lesione insieme a informazioni del caso; il sistema elabora l'immagine in una pipeline a piu' stadi e arriva a una decisione di triage.

\subsection{Evidence E2 - fallback senza immagine}
\textbf{Source:} \texttt{OR2\_Architecture\_Document.pdf}, sezione 2, pagina 2.

\begin{quote}
\small ``Fallback: When no image is available, symptom-only scoring derives urgency from B4 chatbot interaction fields [...]''
\end{quote}

\textbf{Source proposition minima:} se l'immagine non e' disponibile, DermaTriage continua a derivare l'urgenza usando informazioni sintomatologiche del caso.

\begin{ddtaworking}[title={Cosa NON viene promosso da questa evidence}]
Il nome B4, i singoli campi sintomatologici e il dettaglio dello scoring non vengono automaticamente promossi a Decision. Devono essere classificati separatamente come boundary, FR semantics o realization evidence quando arriveremo al livello appropriato.
\end{ddtaworking}

\section{Candidate commitment e semantic owner}

\subsection{Candidate Decision}
\textbf{DEC-01 - Continuita' del triage in assenza di immagine}

\textbf{Context.} Il caso dermatologico puo' essere disponibile con un'immagine della lesione oppure senza immagine. La documentazione descrive un percorso principale image-based e un fallback symptom-only.

\textbf{Decision.} DermaTriage deve consentire la valutazione di triage anche quando non e' disponibile un'immagine della lesione. In tale condizione, la determinazione dell'urgenza deve basarsi sulle informazioni sintomatologiche disponibili sul caso.

\textbf{Consequences.} La valutazione di triage deve supportare almeno le due condizioni ``con immagine'' e ``senza immagine''. La documentazione non autorizza a considerare automaticamente equivalenti gli output dei due percorsi oltre alla derivazione dell'urgenza.

\subsection{Perche' il semantic owner candidato e' Decision}
Il significato non e' soltanto ``eseguire uno scoring symptom-only'': quello sarebbe gia' vicino a un obbligo funzionale. Il commitment piu' stabile e' che \textbf{l'immagine non costituisce una precondizione assoluta per poter effettuare il triage}. Questa scelta restringe MR-01 senza descrivere ancora il comportamento operativo completo.

\section{Decision gate - domande esplicite}

\subsection{Q-D2.1 - Esiste un material project commitment?}
\begin{ddtacheck}
\textbf{Domanda.} La fonte governa una posizione progettuale materialmente significativa, oppure descrive soltanto un dettaglio esecutivo?
\end{ddtacheck}

\textbf{Risposta.} SI. La possibilita' di effettuare il triage in assenza di immagine modifica le condizioni di applicabilita' della capability. Una soluzione alternativa avrebbe potuto richiedere l'immagine e rifiutare o rinviare il caso in sua assenza.

\textbf{Esito: PASS.}

\subsection{Q-D2.2 - La Decision restringe esattamente un MR?}
\begin{ddtacheck}
\textbf{Domanda.} DEC-01 restringe un solo MacroRequirement, oppure introduce responsabilita' autonome appartenenti ad altri branch?
\end{ddtacheck}

\textbf{Risposta.} Restringe \texttt{MR-01}. Stabilisce come la responsabilita' di valutazione di triage si comporta rispetto alla disponibilita' dell'immagine; non introduce una macro-responsabilita' separata.

\textbf{Esito: PASS.}

\subsection{Q-D2.3 - E' una tautologia dell'MR?}
\begin{ddtacheck}
\textbf{Domanda.} Se rimuoviamo DEC-01, MR-01 implica gia' necessariamente che il triage debba funzionare senza immagine?
\end{ddtacheck}

\textbf{Risposta.} NO. MR-01 governa la valutazione di triage dalle informazioni disponibili, ma non stabilisce da solo se l'assenza dell'immagine renda il caso non processabile.

\textbf{Esito: PASS - non tautologica.}

\subsection{Q-D2.4 - Context, Decision e Consequences sono materiali?}
\begin{ddtacheck}
\textbf{Domanda.} Il Context espone una tensione reale, la Decision dichiara una posizione e le Consequences producono effetti downstream materiali?
\end{ddtacheck}

\textbf{Risposta.} SI. La tensione e' la disponibilita' o meno dell'immagine; la posizione e' continuare il triage; le conseguenze riguardano input ammessi, branch funzionali e output che potranno essere richiesti downstream.

\textbf{Esito: PASS.}

\section{Split e non-ridondanza}

\subsection{Q-SPLIT.1 - Con immagine e senza immagine sono Decision indipendenti?}
\begin{ddtacheck}
\textbf{Domanda.} Le due condizioni possono cambiare indipendentemente al punto da richiedere due commitment autonomi con trade-off separati?
\end{ddtacheck}

\textbf{Risposta.} Non emerge evidence sufficiente per due Decision. Il significato governato qui e' uno: la continuita' del triage quando l'immagine manca. Il percorso image-based resta la condizione principale documentata e verra' ulteriormente analizzato attraverso altri commitment/FR.

\textbf{Esito: PASS - NO SPLIT per DEC-01.}

\section{Responsibility-boundary gate}

\subsection{Q-RB.1 - DermaTriage possiede la capability o consuma un servizio esterno?}
\begin{ddtacheck}
\textbf{Domanda.} La responsabilita' di produrre la valutazione di triage e' del progetto DermaTriage oppure appartiene a B4/chatbot come servizio esterno?
\end{ddtacheck}

\textbf{Evidence.} La sezione API distingue gli endpoint DermaTriage dalle ``B4 External API Calls''. B4 fornisce consulti, documenti e campi di interazione e riceve risultati/validazioni; la pipeline di triage e' descritta come responsabilita' DermaTriage.

\textbf{Risposta.} DermaTriage possiede la responsabilita' di triage; B4 e' una capability/integrazione consumata per ottenere o persistere informazioni del caso.

\textbf{Esito: PASS.}

\begin{ddtaworking}[title={Boundary preserved}]
Non trasformare ``B4 chatbot'' nel subject normativo di DEC-01. Il progetto puo' cambiare il mezzo di acquisizione delle informazioni sintomatologiche senza cambiare il commitment fondamentale della Decision.
\end{ddtaworking}

\section{Decision-vs-current-realization test}

\subsection{Q-REAL.1 - Neutralizzando i nomi tecnici resta una scelta?}
\begin{ddtacheck}
\textbf{Domanda.} Rimuovendo B4, chatbot e dettagli dello scoring, resta un commitment governato leggibile?
\end{ddtacheck}

\textbf{Risposta.} SI: ``il triage deve poter continuare anche senza immagine, usando le informazioni sintomatologiche disponibili''.

\textbf{Esito: PASS.}

\subsection{Q-REAL.2 - Cambiare la realizzazione preservando MR-01 cambierebbe la strategia?}
\begin{ddtacheck}
\textbf{Domanda.} Possiamo sostituire B4 o il meccanismo di scoring mantenendo invariata DEC-01?
\end{ddtacheck}

\textbf{Risposta.} SI. Questo dimostra che B4 e lo scoring corrente non sono l'identita' della Decision. Al contrario, passare a una policy ``no image -> no triage'' cambierebbe il commitment pur lasciando invariato l'obiettivo macro di MR-01.

\textbf{Esito: PASS - Decision separata dalla realization.}

\section{Decision kind review}

\subsection{Q-KIND.1 - SELECTABLE, NECESSITY-CONSTRAINED o DEFAULT?}
\begin{ddtacheck}
\textbf{Domanda.} Esiste almeno un'alternativa materialmente distinta concepibile oppure la posizione e' una necessita' inevitabile del dominio?
\end{ddtacheck}

\textbf{Risposta.} Esiste una alternativa materialmente distinta: richiedere obbligatoriamente l'immagine e non effettuare il triage quando manca. La fonte documenta la posizione scelta, ma non documenta il processo con cui l'alternativa e' stata valutata.

\textbf{Kind candidate: SELECTABLE.}

\begin{tabularx}{\textwidth}{L{55mm}Y}
\toprule
\textbf{Proposition} & \textbf{Evidence status} \\
\midrule
Il sistema continua il triage senza immagine & AFFIRMED \\
Esiste una alternativa materialmente distinta concepibile & AFFIRMED per review strutturale; non equivale a storia decisionale documentata \\
Il team ha formalmente valutato quella alternativa & NOT SPECIFIED \\
Rationale storico della scelta & NOT SPECIFIED \\
\bottomrule
\end{tabularx}

\section{Semantic-sufficiency check}

\subsection{Q-SEM.1 - Esistono letture materialmente diverse?}
\begin{ddtacheck}
\textbf{Domanda.} Due reviewer potrebbero leggere ``fallback symptom-only'' in modi che producono downstream structure diversa?
\end{ddtacheck}

\textbf{Risposta.} SI. Una lettura potrebbe assumere che il fallback produca l'intero outcome del percorso image-based; un'altra che produca soltanto l'urgenza.

\subsection{Q-SEM.2 - Smallest critical proposition}
\begin{ddtacheck}
\textbf{Domanda.} Qual e' la proposizione minima che separa le due letture?
\end{ddtacheck}

\textbf{Smallest critical proposition:} ``Il percorso symptom-only deve produrre l'intero insieme di output del percorso image-based, oppure soltanto l'urgenza?''

\subsection{Q-SEM.3 - La fonte governa la proposizione?}
La fonte afferma esplicitamente soltanto che il fallback symptom-only \emph{derives urgency}. Non afferma equivalenza di pathology prediction, reasoning, clinical description, similar cases o altri output.

\textbf{Evidence status: NOT SPECIFIED oltre all'urgenza.}

\subsection{Q-SEM.4 - Semantic owner corretto}
Il gap non richiede di ampliare DEC-01. La Decision deve governare la continuita' del triage; la forma e cardinalita' degli output appartiene alla successiva review funzionale, se e quando la fonte la governa.

\textbf{Semantic sufficiency disposition: PASS con gap preservato.}

\section{Disposition e STOP}

\begin{ddtacore}[title={DEC-01 disposition}]
\textbf{DEC-01: PASS}\par
\textbf{Kind: SELECTABLE}\par
\textbf{Parent: MR-01}\par
\textbf{NS-01:} rationale storico e alternative formalmente considerate - NOT SPECIFIED.\par
\textbf{NS-02:} equivalenza dell'output completo symptom-only rispetto al percorso image-based - NOT SPECIFIED; e' affermata soltanto la derivazione dell'urgenza.
\end{ddtacore}

\subsection{STOP question}
\begin{ddtacheck}
\textbf{Domanda.} Dopo DEC-01, esiste ancora significato governato che restringe materialmente MR-01 e che non puo' essere ridotto a FR/realization di DEC-01?
\end{ddtacheck}

\textbf{Risposta.} SI. La documentazione distingue almeno un risultato analitico di urgenza \texttt{HIGH/MEDIUM/LOW} e una successiva assegnazione operativa \texttt{P1-P4} con SLA. Questa relazione deve essere riesaminata come possibile commitment autonomo senza inglobarla in DEC-01.

\textbf{STOP: NO - continuare la Decision review di MR-01.}

\section{Osservazione metodologica aperta}

\begin{ddtaworking}[title={OBS-DDTA-DERMA-01 - Gate visibility}]
\textbf{Osservazione:} una tabella che riporta soltanto ``Gate / PASS / motivazione'' puo' essere troppo compressa per un artifact didattico o di ricerca, anche se il reviewer ha effettivamente applicato i test.

\textbf{Ipotesi da verificare:} il working review record potrebbe richiedere per ogni gate almeno \emph{question + answer + evidence + disposition}, mentre la project documentation finale deve restare priva di questo commentary.

\textbf{Stato: OPEN OBSERVATION - nessuna modifica alla metodologia R4 in questa fase.}
\end{ddtaworking}

\section{Formato vincolante per la prossima Decision}

Per DEC-02 la review verra' esposta nella stessa sequenza, senza saltare le domande:

\begin{lstlisting}
1. Source evidence
2. Candidate commitment
3. Semantic owner
4. Decision gate questions + answers
5. Split / non-redundancy questions
6. Responsibility-boundary question
7. Decision-vs-realization questions
8. Decision kind question
9. Semantic-sufficiency questions
10. Explicit disposition
11. STOP question
\end{lstlisting}

Nessun FunctionalRequirement viene creato fino alla chiusura della Decision review pertinente.

\section{Valutazione metodologica della sessione}

Questa sezione chiude la nota di lettura con una valutazione della metodologia \emph{limitata a quanto osservato fino a DEC-01}. Non e' ancora la valutazione conclusiva del holdout DermaTriage: pregi, difetti e proposte restano soggetti a conferma sulle Decision successive, sugli FR e sul documentation gate finale.

\subsection{Pregi osservati}

\begin{tabularx}{\textwidth}{L{42mm}Y L{25mm}}
\toprule
\textbf{Pregio} & \textbf{Evidenza emersa in DEC-01} & \textbf{Stato} \\
\midrule
Separazione tra significato governato e realizzazione & Dal termine sorgente ``fallback'' e' stato estratto il commitment stabile ``il triage puo' continuare senza immagine'', senza trasformare automaticamente B4, chatbot o scoring in Decision. & CONFERMATO \\
Preservazione esplicita dell'incertezza & Il metodo ha impedito di assumere che il percorso symptom-only produca pathology, reasoning o tutti gli output del percorso image-based: tali aspetti restano \texttt{NOT SPECIFIED}. & CONFERMATO \\
Controllo del semantic owner & La continuita' del triage resta nella Decision; input sintomatologici, output e comportamento operativo vengono rinviati al livello funzionale appropriato invece di essere accumulati nella stessa entita'. & CONFERMATO \\
Responsibility-boundary discipline & La review distingue la responsabilita' di DermaTriage dalla capability esterna B4, evitando di attribuire al progetto ownership di servizi consumati. & CONFERMATO \\
Stopping discipline & Il PASS di DEC-01 non chiude artificialmente MR-01: la STOP question rende esplicito che esistono ancora commitment materiali da riesaminare. & CONFERMATO \\
Tracciabilita' del ragionamento & Quando le domande sono mostrate esplicitamente, il lettore puo' ricostruire perche' una proposizione e' diventata Decision, gap o realization evidence. & CONFERMATO, con condizione di visibilita' \\
\bottomrule
\end{tabularx}

\subsection{Difetti, costi e rischi osservati}

\begin{tabularx}{\textwidth}{L{42mm}Y L{25mm}}
\toprule
\textbf{Costo / difetto} & \textbf{Manifestazione nella sessione} & \textbf{Stato} \\
\midrule
Elevato carico interpretativo sul reviewer & Distinguere Decision da FR e da realization non deriva meccanicamente dalla parola ``fallback''; richiede neutralizzazione dei nomi tecnici, confronto di alternative e ownership review. & REALE \\
Rischio di compressione eccessiva dei gate & La prima esposizione di DEC-01 riportava gli esiti ma non rendeva visibili tutte le domande applicate; il metodo puo' risultare meno auditabile se il working record comprime troppo il processo. & REALE \\
Verbosity procedurale & Esplicitare domanda, evidence, risposta ed esito per ogni gate rende la review piu' lunga della documentazione progettuale ordinaria. Il costo e' accettabile nel working record, ma sarebbe eccessivo nel documento finale. & REALE \\
Possibile ambiguita' tra alternativa concepibile e alternativa storicamente valutata & La classificazione \texttt{SELECTABLE} richiede riconoscere che esiste una posizione materialmente diversa, senza inventare che il team l'abbia effettivamente considerata. Questa distinzione richiede attenzione. & REALE \\
Rischio di layer forcing & La gerarchia puo' spingere il reviewer a cercare una Decision anche quando la fonte contiene soltanto un obbligo funzionale. Il counterexample gate della R4 mitiga il rischio, ma dovra' essere verificato su altri casi. & DA VERIFICARE \\
\bottomrule
\end{tabularx}

\subsection{Miglioramenti metodologici candidati}

\begin{ddtaworking}[title={Miglioramenti da verificare, non ancora promossi}]
\begin{enumerate}
  \item \textbf{Gate visibility contract.} Nel working review record rendere obbligatorio, per ogni gate materialmente applicabile, il formato minimo: \texttt{question -> source evidence -> answer -> disposition}. Il documento finale governato rimane invece sintetico.
  \item \textbf{Evidence pointer esplicito.} Collegare ogni risposta critica al frammento o alla posizione della fonte che la sostiene, distinguendo evidence diretta da inferenza di review.
  \item \textbf{Selectable-kind clarification.} Specificare nella guida che ``esiste una alternativa materialmente distinta'' non significa ``la fonte documenta che il team l'ha valutata''. La seconda proposizione deve restare \texttt{NOT SPECIFIED} se manca evidence.
  \item \textbf{Decision extraction worksheet.} Fornire una scheda didattica riusabile con le domande gia' in sequenza, cosi' il reviewer non puo' saltare accidentalmente responsibility boundary o Decision-vs-realization.
\end{enumerate}
\end{ddtaworking}

\subsection{Valutazione complessiva dopo DEC-01}

\begin{ddtacore}[title={Valutazione provvisoria}]
Su DEC-01 la metodologia sta migliorando la documentazione DermaTriage in modo osservabile: trasforma una descrizione architetturale centrata sulla soluzione in un commitment progettuale piu' stabile, conserva i gap invece di completarli per inferenza e separa ownership, comportamento e realization.

Il principale costo osservato non e' una perdita di significato, ma \textbf{l'onere di review}: la qualita' dipende dalla capacita' di applicare i gate in modo disciplinato e di renderne visibile il ragionamento. La prima esecuzione compressa di DEC-01 mostra che questo rischio e' concreto anche quando le regole esistono gia' nella guida.

\textbf{Conclusione della sessione:} nessuna modifica a R4 e' giustificata ancora. \texttt{OBS-DDTA-DERMA-01} resta aperta e viene verificata sulle prossime Decision. Se la necessita' di esplicitare le domande ricorre sistematicamente, il problema diventera' candidato a miglioramento della guida e non semplice errore di esecuzione.
\end{ddtacore}

\end{document}
